\subsection{Tilted shells and exact exterior cuts}
\label{sec:exact-shell}

In this subsection, $R=\sqrt6$ is the outer radius of the shell.
All other centres of an extension have norm at least $R$.
The bounds concern the actual Voronoi cell $V_0$.

\begin{thm}[Angular stability for nearly equal radii]\label{thm:exact-shell-tilt}
Let $u_i$ be the twenty-four unit directions $(\pm e_j\pm e_k)/\sqrt2$, $j<k$, of $D_4$.
Suppose that $Y=\{2a_iw_i:1\le i\le24\}$, where $w_i$ are unit vectors, $1\le
a_i\le\sqrt{3/2}$, and the matching to the $u_i$ has been made after one common rotation.
Set

\[
\eta=\max_i(a_i-1),\qquad m=\min_i a_i,\qquad h=1+\eta,\qquad
\epsilon=\max_i|w_i-u_i|.
\]

The following bounds hold for the actual cell.

\begin{enumerate}[label=(\roman*)]
\item If all $a_i=h$ and $\epsilon\le\eta/2$, then $\vol V_0\ge8+\eta/8$.
\item If $m\ge1+9\eta/10$ and $\epsilon\le2\eta/5$, then $\vol V_0\ge8+\eta$.
\end{enumerate}
At $\eta=0$ the cell is the regular twenty-four-cell.
\end{thm}
\begin{proof}
Let $Q=\{x:x\cdot u_i\le1\text{ for all }i\}$ and $K_0=\operatorname{conv}\{2u_i\}$. In
coordinates $Q$ is given by $|x_j|+|x_k|\le\sqrt2$. In the region where $x_1=t\ge0$ is a
largest coordinate in absolute value, the other three coordinates independently have
absolute value at most $\min(t,\sqrt2-t)$. Thus

\[
\vol Q=8\int_0^{\sqrt2}8\min(t,\sqrt2-t)^3\,dt=8.
\]

The vertices of $Q$ are $\sqrt2 g$ with $g\in G=\{\pm
e_j\}\cup\{(\pm1,\pm1,\pm1,\pm1)/2\}$. Indeed, the portion with every coordinate of
absolute value at most $1/\sqrt2$ is a cube. In the region $t=x_1\ge1/\sqrt2$ with $x_1$
largest, the transverse cube of side $2(\sqrt2-t)$ interpolates linearly between a facet
of that cube and $\sqrt2e_1$. Coordinate signs and permutations give all other portions.
In particular $Q\subset B(\sqrt2)$. By polarity $K_0=2Q^\circ$, so its facet inequalities
are $x\cdot g\le\sqrt2$ for $g\in G$, and $B(\sqrt2)\subset K_0$.

For $y=2aw\in Y$, the point $2k(a)w$, where $k(a)=3a/(2a^2+1)$, is on the origin's side of
every exterior bisector. In fact $2k(a)w=R^2y/(R^2+|y|^2-4)$. If $z$ is exterior, then
$2y\cdot z\le |y|^2+|z|^2-4$, which implies

\[
2\frac{R^2y\cdot z}{R^2+|y|^2-4}\le |z|^2
\]

by $(|y|^2-4)(|z|^2-R^2)\ge0$. The ball $B(R/2)$ is protected as well. Their convex hull
$K_R(Y)$ is therefore protected.

Put $b=\sqrt2\epsilon$ and $s=m/(1+b)$. For $x\in Q$, $x\cdot w_i\le1+b$, so $sQ\subset
V(Y)$. The function $k$ decreases for $a\ge1$, and $K_R(Y)$ contains $0$, so it contains
each $2k(h)w_i$. The support functions of the hulls of $2u_i$ and $2w_i$ differ by at most
$2\epsilon|v|$ in direction $v$. Since the former contains $B(\sqrt2)$, the latter
contains $(1-b)K_0$ whenever $b<1$. Hence $sQ\cap k(h)(1-b)K_0\subset V_0$.

The portion of $sQ$ beyond one facet of $k'K_0$, $k'=k(h)(1-b)$, has volume at most
$8(s-k')_+^4$, where $t_+=\max(t,0)$. For the facet with normal $e_1$, at height
$x_1\ge\sqrt2k'$ the three other coordinates lie in a cube of side $2(\sqrt2s-x_1)$. Its
integral gives the bound. Signed permutations and the orthogonal Hadamard matrix with rows
$(1,1,1,1),(1,1,-1,-1),(1,-1,1,-1),(1,-1,-1,1)$, divided by $2$, preserve the roots and
carry this normal to all of $G$. Therefore

\[
\vol V_0\ge 8s^4-192(s-k')_+^4.
\]

The angular terms in the width of the removed cap have opposite signs. More precisely,

\[
s-k'=(m-k(h))(1-b)+\frac{mb^2}{1+b}.
\]

In both families $mb/(1+b)\le m-1\le m-k(h)$. To check the first inequality, use $m<5/4$,
$b\le71\eta/100$ in the equal-radius case, and $b\le57\eta/100$ in the other. Thus
$s-k'\le m-k(h)\le h-k(h)\le(11/8)\eta$. For the last inequality, when $h>1$, multiply
$(h-k(h))/(h-1)=2h(h+1)/(2h^2+1)$ by its positive denominator and use $6h^2-16h+11>0$. It
also holds at $h=1$.

Since $0\le\eta<9/40$, the cap loss is at most $(32019867/4096000)\eta$, using
$192(11/8)^4(9/40)^3=32019867/4096000$. In the equal-radius case
$32(s-1)\ge(37120/4639)\eta$. In the other case $32(s-1)\ge(42240/4513)\eta$. These follow
by bounding $s-1=(m-1-b)/(1+b)$ with the respective bounds on $b$ and $\eta$. Both have
$s\ge1$, so $8(s^4-1)\ge32(s-1)$. Subtracting the cap loss leaves more than $\eta/8$ and
$\eta$, respectively, when $\eta>0$. At $\eta=0$, $V(Y)=Q\subset K_0\subset K_R(Y)$. This
proves the bounds.
\end{proof}

\begin{example}[A tilted shell with $T(Y)<8$]\label{ex:exact-shell-tilt}
Set $\eta=1/1000$ and $a_i=1+\eta$ for all $i$. Choose two roots at angle $\pi/3$. Rotate
each towards the other in their common plane by the angle $\theta$ of chord length
$e=\eta/2$. Leave the other roots fixed. The inner product of the two tilted directions is

\[
\cos(\pi/3-2\theta)=\tfrac12+\tfrac{\sqrt3}{2}\sin(2\theta)-\sin^2\theta.
\]

Its excess above $1/2$ is between $4\eta/5$ and $9\eta/10$. Indeed,
$\sin\theta=e\sqrt{1-e^2/4}\le e$, $\cos\theta=1-e^2/2$, and
$(98/100)\eta<\sin(2\theta)\le\eta$. Any other pair has inner product at most $1/2+e$. The
packing threshold $1-1/(2(1+\eta)^2)$ is strictly above $1/2+9\eta/10$, as
$(2+\eta)/(2(1+\eta)^2)>9/10$. This verifies the packing. The earlier small-angle
criterion $\epsilon\le(\sum_i(a_i-1))/100$ cannot hold under any matching, since it would
bound every inner product by $1/2+2(24\eta)/100$.

Put $r_0=R/2$. The twenty-four caps of $Q$ beyond $x\cdot g=r_0$ lie outside $B(r_0)$, and
each has volume $2(\sqrt2-r_0)^4$ by the cube-section integral. They are disjoint up to
boundaries since distinct $g,g'\in G$ have $g\cdot g'\le1/2$ and $x\in Q$ has
$|x|\le\sqrt2$, giving $x\cdot(g+g')\le2r_0$. For $x\in V(Y)$, $\max_i x\cdot
u_i\le1+\eta+e|x|$ and $|x|\le\sqrt2\max_i x\cdot u_i$. Hence $V(Y)\subset t'Q$ for
$t'=(1+\eta)/(1-\sqrt2e)$. Since $t'\ge1$, $V(Y)\cap B(r_0)\subset t'(Q\cap B(r_0))$.
Using $\sqrt2-r_0>189/1000$ and $t'<4004/3997$ gives the rational bound

\[
\vol(V(Y)\cap B(r_0))
< \left(\frac{4004}{3997}\right)^4\left(8-48\left(\frac{189}{1000}\right)^4\right)
<8.
\]

Thus the truncated cell falls below the target in this example. The new estimate gives
$\vol V_0\ge8+1/8000$ for every extension with exterior centres at norm at least $R$.
\end{example}

\begin{prop}[Protection at the vertices of the regular cell]\label{prop:exact-shell-vertices}
Use the same matching, with $\epsilon\le1/32$. For a vertex $v=\sqrt2g$ of $Q$, let $h_v$
and $\epsilon_v$ be the maxima of $a_i$ and $|w_i-u_i|$ over the six roots satisfying
$u_i\cdot g=1/\sqrt2$. Define $q_v=1/\sqrt2-\epsilon_v$ and

\[
\lambda_v=\frac{\max\bigl(R,\ 2h_vq_v+2\sqrt{1-h_v^2(1-q_v^2)}\bigr)}{2\sqrt2}.
\]

Every exterior bisector leaves $\lambda_v v$ on the origin's side. Therefore $V_0$
contains $V(Y)$ intersected with the convex hull of $K_R(Y)$ and all twenty-four points
$\lambda_v v$.
\end{prop}

\begin{proof}
Use a root symmetry to take $g=e_1$. If an exterior centre $z=su$, $|u|=1$, cuts
$\lambda_v v$, then $q=u\cdot e_1>s/(2\sqrt2\lambda_v)\ge\sqrt3/2$. Here $\lambda_v\le1$
follows because $a\mapsto 2aq_v+2\sqrt{1-a^2(1-q_v^2)}$ is nonincreasing for $1\le
a\le\sqrt{3/2}$ and at $a=1$ has value $4q_v\le2\sqrt2$. For the monotonicity,
differentiating and squaring the nonnegative terms reduces the sign to $q_v^2\le
a^2(1-q_v^2)$. The square root is real since $q_v>2/3$ and $a^2\le3/2$.

Among the six roots $(e_1\pm e_j)/\sqrt2$ there is one with

\[
u\cdot u_i\ge \frac{q+\sqrt{(1-q^2)/3}}{\sqrt2}\ge\frac1{\sqrt2}.
\]

The second inequality follows by squaring $\sqrt{(1-q^2)/3}\ge1-q$, valid for $1/2\le
q\le1$. For this root $u\cdot w_i\ge q_v$. Separation of $z$ from $2a_iw_i$ gives
$s^2-4a_iq_vs+4(a_i^2-1)\ge0$. The smaller root is below $R$ since it is at most
$2a_iq_v\le\sqrt3$. The larger root is at least its value at $a_i=h_v$ by monotonicity.
Thus $s\ge2\sqrt2\lambda_v$, contradicting $s<2\sqrt2\lambda_v q\le2\sqrt2\lambda_v$.
Convexity gives the asserted protected hull.
\end{proof}

\begin{remark}[Sharp scale for an equal radial displacement]
When the directions are exactly the roots and all $a_i=a$, this protection is sharp for
homothetic copies of $Q$. The scale is $t=(a+\sqrt{2-a^2})/2$. The possible exterior
centre $2\sqrt2t\,e_1$ lies at distance at least $2$ from every $2au_i$, with equality for
the six roots at the vertex $\sqrt2e_1$, and its bisector cuts any larger homothetic copy
of $Q$. The radial loss is

\[
1-t=\frac{(a-1)^2}{2-a+\sqrt{2-a^2}}.
\]

For comparison, $K_R(Y)$ has support $\sqrt2k(a)$ in the direction $e_1$. This support
exceeds $R/2$ and the deficit $1-k(a)=(a-1)(2a-1)/(2a^2+1)\ge(a-1)/4$ is larger than $1-t$
for $1<a\le\sqrt{3/2}$. Thus the larger protected region recovers positive volume even
after intersecting with $V(Y)=aQ$. In fact it recovers at least $192(t-k(a))^4$ from the
twenty-four caps of $tQ$ beyond the planes $x\cdot g=\sqrt2k(a)$. In each cap the
transverse section is the full cube because $k(a)>t/2$, so the integral is exact. Two caps
cannot overlap since distinct $g,g'\in G$ have $g\cdot g'\le1/2$, $|x|\le\sqrt2t$ on $tQ$,
and $k(a)>(\sqrt3/2)t$.
\end{remark}

\begin{prop}[Finite exterior-cut test]\label{prop:exact-shell-test}
Let $Y$ be a shell configuration, with arbitrary directions. Define the compact set

\[
F_Y=\{p:|p|^2\le4/R^2,\quad (|y|^2-4)|p|^2-4y\cdot p+4\ge0\text{ for all }y\in Y\}.
\]

Then the common region of the actual cells over all allowed extensions is

\[
V(Y)\cap\{x:\max_{p\in F_Y}x\cdot p\le1\}.
\]
The support maximum can be computed from subsets of at most four boundary equations.
\end{prop}
\begin{proof}
Invert a possible exterior centre by $p=2z/|z|^2$. The inequalities $|z|\ge R$ and
$|z-y|\ge2$ become exactly the defining inequalities of $F_Y$ for $p\ne0$, and the
corresponding bisector is $x\cdot p=1$. The point $p=0$ adds no constraint. Every feasible
$p\ne0$ can be realised by adjoining the single centre $z=2p/|p|^2$.

Consider subsets of at most four boundary equations. Each boundary equation defines a
sphere or a hyperplane. An intersection of such equations is affine if there are only
hyperplanes. Otherwise, subtract the equations of the spheres after normalising the
coefficients of $|p|^2$, to obtain one sphere intersected with an affine space. Discard
empty intersections. For an isolated point or two-point intersection, test all the points.
For a sphere of positive dimension in its affine space, choose the point maximising the
first linear functional in the list $x\cdot p,e_1\cdot p,\ldots,e_4\cdot p$ that is
nonconstant on that sphere. It is given by the centre plus the radius times the normalised
orthogonal projection of that functional's coefficient vector onto the direction space.
Check each candidate against all constraints defining $F_Y$. Positive-dimensional purely
affine intersections give no candidates. Take the largest value of $x\cdot p$ among the
feasible candidates, also allowing $p=0$.

Take a lexicographic maximiser for that list on $F_Y$, which exists by compactness. Near
it, all constraints except the active equations are strict. A positive-dimensional affine
intersection of the active equations cannot have this maximiser, since the first
functional nonconstant on the affine space has a local increasing direction. On a sphere
of positive dimension the lexicographic maximiser must locally maximise the first
nonconstant functional, hence is its unique maximiser on the sphere. The intersection of
the full active set is defined by at most four equations if there are only planes. With a
sphere, retain one sphere and a basis of the linear equations obtained by subtracting it
from the other normalised sphere equations and adjoining the plane equations. If this
basis has rank at most three, at most four original equations suffice. If it has rank
four, keep just three independent linear equations with the sphere. The resulting
intersection is a line intersected with a sphere, so it has at most two points and
includes the maximiser. This proves exhaustiveness in every case. With rational $Y$, $R^2$
and $x$, the operations on a given subset require rational arithmetic and at most a square
root for each candidate, and feasibility comparisons can be made exactly.
\end{proof}
\subsection{Stability and exact computations for the shell test}
\label{sec:shell-computation}

Write \(C_R(Y)\) for the common region in \cref{prop:exact-shell-test},
including the half-spaces of \(Y\).

\begin{lem}[Stability under perturbation of a shell]
\label{lem:shell-certificate-stability}
Let \(Y=\{y_1,\ldots,y_n\}\) and \(Y'=\{y'_1,\ldots,y'_n\}\)
be two shell configurations in \(2\le |y|\le R=\sqrt6\), with
\(\max_i|y_i-y_i'|\le\delta<1\). If \(P\subset C_R(Y)\cap B(M)\), then
\[
 \lambda P\subset C_R(Y'),\qquad
 \lambda=\min\left(1-\delta,\frac{1}{1+(M+1)\delta/2}\right).
\]
In particular, a polytope \(P\) of volume \(L>8\) gives
\(\vol V_0'\ge\lambda^4 L>8\) on any neighbourhood for which
\(\lambda^4 L>8\).
\end{lem}

\begin{proof}
For each unit direction \(u\), let \(s_y(u)\) be the least \(s\ge R\)
satisfying \(|su-y|\ge2\). The squared distance is nondecreasing in \(s\)
on \([R,\infty)\) since \(u\cdot y\le R\). Thus
\(s_Y(u)=\max_{y\in Y}s_y(u)\), with value \(R\) for an empty shell.
The formula of \cref{shell-protect:prop} expresses the exterior part of
\(C_R(Y)\) as \(x\cdot u\le s_Y(u)/2\) for every \(u\).

On the segment \(y(t)=(1-t)y+ty'\), \(0\le t\le1\), we have \(|y(t)|\le R\).
Where \(s=s_{y(t)}(u)>R\), differentiation of \(|su-y(t)|^2=4\) gives
\[
 \left|\frac{ds}{dt}\right|
 \le \frac{2|y'-y|}{s-u\cdot y(t)},\qquad
 s-u\cdot y(t)=\frac{s^2+4-|y(t)|^2}{2s}\ge\frac{2}{R}.
\]
The last inequality uses the monotonicity of \((s^2+4-R^2)/(2s)\) for
\(s\ge R=\sqrt6\). On the constant branch \(s=R\) the derivative is zero.
The branches meet continuously. The transition points lie among the
roots of the quadratic \(|Ru-y(t)|^2=4\), with the constant branch
throughout if that equation holds identically. The derivative bound
therefore gives \(|s_y(u)-s_{y'}(u)|\le R|y-y'|\). Taking maxima yields
\(s_{Y'}(u)\ge s_Y(u)-R\delta\ge(1-\delta)s_Y(u)\).
Since \(s_Y(u)>0\), scaling by any \(0<\lambda\le1-\delta\) preserves
the exterior constraints with the new right-hand sides.

For the constraints from \(Y'\), write \(y_i=r_i u_i\) and
\(y_i'=r_i' u_i'\) with unit directions. Then
\[
 |r_i-r_i'|\le\delta,\qquad
 |y_i-y_i'|^2=(r_i-r_i')^2+r_i r_i'|u_i-u_i'|^2.
\]
In particular \(|u_i-u_i'|\le\delta/2\).
For \(x\in P\) we get
\[
 x\cdot u_i'\le\frac{r_i}{2}+\frac{M\delta}{2}
 \le\frac{r_i'}2+\frac{(M+1)\delta}2
 \le\frac{r_i'}2\left(1+\frac{(M+1)\delta}{2}\right).
\]
Scaling by \(\lambda\) proves the new interior inequalities. The
inclusion and volume bound follow.
\end{proof}

\begin{example}[A certified neighbourhood of a rational shell]
\label{ex:exact-shell-computed}
Use the orthogonal image of the \(D_4\) directions
\[
 U=\{\pm e_i:1\le i\le4\}
   \ \cup\ \{(\pm1,\pm1,\pm1,\pm1)/2\}.
\]
These coordinates result from applying the two-dimensional Hadamard
transform on coordinate pairs \((1,2)\) and \((3,4)\). Put
\(a=1001/1000\), \(c=99999/100000\) and \(Y=2aU\).
The uncut cell \(V(Y)=aQ\) is given by \(|x_i|\le a\) and
\(\sum_i|x_i|\le2a\). Consider its intersection \(P\) with the
twenty-four half-spaces \(\pm x_i\pm x_j\le2c\), \(i<j\).
They truncate the vertices of \(aQ\), which have two coordinates
equal to \(\pm a\) and two equal to zero. The cuts are disjoint on
\(aQ\) because \(c/a>\sqrt3/2\). Indeed the outward unit directions
of distinct cuts have inner product at most \(1/2\), and \(aQ\) has
circumradius \(\sqrt2a\).

The vertices of \(P\) are the intersections of the cutting planes
with the edges incident to the truncated vertices. In the region
cut off by \(x_1+x_2=2c\), both \(x_1,x_2\) are positive and
at most \(a\). Setting \(b_1=a-x_1\), \(b_2=a-x_2\), the region is
\(b_1,b_2\ge0\), \(b_1+b_2<2(a-c)\), and
\(|x_3|+|x_4|\le b_1+b_2\). The truncated vertices in that plane
therefore have \((b_1,b_2)=(0,2(a-c))\) or its swap, and
\((x_3,x_4)\) a vertex of the transverse diamond. Signed coordinate
permutations preserve \(Y\) and act transitively on these points
over all the planes. One representative is
\[
 (a,\,2c-a,\,2(a-c),\,0).
\]
The exact test verifies protection of this point, and finds a cutting
centre for the untrimmed vertex \((a,a,0,0)\). Convexity gives
\(P\subset C_R(Y)\). Integrating the diamond of area \(2(b_1+b_2)^2\)
over the nonnegative triangle in \((b_1,b_2)\) gives \(8(a-c)^4\)
for each removed piece. Consequently
\[
 \vol P = 8a^4-192(a-c)^4 > 8.032 .
\]
Since \(P\subset B(2)\), \cref{lem:shell-certificate-stability} also gives
\(\vol V_0'>8.007\) for every admissible shell obtained by moving each
centre by at most \(1/2000\). The code checks these rational comparisons
and the two exterior tests.
\end{example}

\begin{example}[Tests at the other three shell counts]
For \(n=25,26,27\), take the twelve equatorial centres
\(\tfrac{17}{10}(\pm e_i\pm e_j)\) for \(1\le i<j\le3\), all eight
points \((\pm1,\pm1,\pm1,17/10)\), and the first \(n-20\) points of
\((\pm1,\pm1,\pm1,-17/10)\) in lexicographic order. All lie in the
shell and satisfy the packing constraints. For each set, the exact
test certifies protection of \((0,0,0,5/4)\) and produces an exterior
cut for \((0,0,0,3/2)\). Both points satisfy the inequalities of the
given shell. The first is outside the standard truncation ball.

Each of these shells has all norms larger than \(12/5+1/500\).
After any admissible movement of at most \(1/500\) per centre,
the actual cell therefore contains \(B(6/5)\), whose volume
\(\tfrac{\pi^2}{2}(6/5)^4\) exceeds \(10.2\). The examples
exercise the exact exterior test at each count. Their radial margin
already gives the volume bound, so they provide no certificate for
the low-volume configurations that remain at these counts.
\end{example}

The supplementary script \texttt{anc/exterior\_test.py} uses rational
arithmetic. For each candidate point it decides signs of
\(A+B\sqrt d\) by comparing \(A^2\) with \(B^2d\) when the two terms have
opposite signs. The checked output is in
\texttt{anc/exterior-computation.log}.
Together with \cref{lem:shell-certificate-stability}, this gives a
way to verify a finite cover of shell configurations by neighbourhoods
with strict volume bounds. Such a cover has yet to be established
for the remaining cases of \cref{prop:reduction-G}.

